

For completeness, we finish with a small number of games that either contain discontinuities or have non-compact strategy sets. Neither equilibria nor best responses necessarily exist in such games; for this reason, we include only examples with known solutions.

\begin{example}
A two-player zero-sum game in $m \times m$ variables on the simplex
\[
x_i \in \Delta^{m-1} = \left\{ z \in \mathbb{R}^m \ \middle|\ z_j \ge 0,\ \sum_{j=1}^m z_j = 1 \right\} \quad \forall i \in \{1,2\}
\]
The utility is defined by the discontinuous function
\[
u_1(x_1, x_2) = \sum_{i=1}^m \operatorname{sign}\left( x_{1j} - x_{2j} \right)
\]
This is an infinite version of the classical \emph{Colonel Blotto} game \cite{Laslier.colonel,Kvasov.colonel,gross.wagner.colonel,Golman2009}, in which all isolated fronts have equal value. Pure strategies exist only for $m \in \{1,2\}$, in which case every strategy is a tie; or for $m=3$, when players should allocate $\frac{1}{3}$ to every front. Otherwise, \textcite{roberson.colonel} proved that players must play a strategy with marginals $\mu^\star_{ij} = \mathcal{U}\left[0,\frac{2}{m}\right]$ for every front $j$.
\end{example}

